% GENERATED FILE. Edit canonical lessons, then run npm run generate:books.
% canonical-source-hash: fnv1a64:f8191bae66384d47
% canonical-lessons: MW-C248-mehdi, MW-C248-parnavno, MW-C248-dayjo, MW-C248-sajno, MW-C248-sehro

\chapter{Henna and Wedding Finery}
\label{ch:mw-henna-and-wedding-finery}
\hlchaptermodality{248}

\begin{chapteropening}
\textbf{By the end of this chapter:} \emph{I can say what happened at a wedding.}

Complete the last lesson of chapter 248: I can say what happened at a wedding.
\end{chapteropening}

\section[\texorpdfstring{\mw{मेंहदी} (mẽhdī)}{मेंहदी (mẽhdī)}]{\mw{मेंहदी} (mẽhdī) — henna}

\label{lesson:MW-C248-mehdi}



\begin{quote}
Before the new one: say the Marwadi for a bride, then the Marwadi for a wedding pavilion.
\end{quote}

\subsection*{You'll want to know: \mw{मेंहदी}}
\textbf{\mw{मेंहदी}} — \emph{mẽhdī} — \textquotedblleft{}henna\textquotedblright{}.

\begin{quote}
\mw{बहन} \mw{हाथ} \mw{माथै} \mw{मेंहदी} \mw{लगावै} \mw{है।}
\end{quote}

\begin{quote}
\emph{bahan hāth māthai mẽhdī lagāvai hai.}
\end{quote}

\begin{quote}
My sister puts henna on her hands.
\end{quote}

A feminine noun.

\begin{cousinweb}
\textbf{\mw{मेंहदी}} (\emph{mẽhdī}) is inherited from Sanskrit \emph{mendhī}, the henna plant. CDIAL gives no Marwari form in this entry; it lists Hindi \emph{mẽhdī} (CDIAL 10328).
\end{cousinweb}

\begin{grammarlens}[title={What you've built}]
One more word for saying what happened.
\end{grammarlens}

\subsection*{Your turn}
\begin{itemize}
\raggedright
  \item \emph{Say it:} \emph{mẽhdī}
  \item \emph{Say it:} \emph{mẽhdī}, once more
  \item \emph{Say it:} \emph{bahan hāth māthai mẽhdī lagāvai hai.} — My sister puts henna on her hands.
  \item \emph{Recall:} say the Marwadi for a bride, then the Marwadi for a wedding pavilion, then say \emph{mẽhdī} again
\end{itemize}

\begin{tcolorbox}[breakable,colback=teal!4,colframe=teal!35!black,title={Before you move on}]
What does \mw{मेंहदी} mean? (\textquotedblleft{}Henna\textquotedblright{}.) Say it once more.
\end{tcolorbox}
\section[\texorpdfstring{\mw{परणावणो} (parṇāvṇo)}{परणावणो (parṇāvṇo)}]{\mw{परणावणो} (parṇāvṇo) — to give in marriage; to marry off}

\label{lesson:MW-C248-parnavno}



\begin{quote}
Before the new one: say the Marwadi for a wedding pavilion, then the Marwadi for henna.
\end{quote}

\subsection*{You'll want to know: \mw{परणावणो}}
\textbf{\mw{परणावणो}} — \emph{parṇāvṇo} — \textquotedblleft{}to give in marriage; to marry off\textquotedblright{}.

\begin{quote}
\mw{म्हैं} \mw{म्हारी} \mw{बेटी} \mw{परणाई।}
\end{quote}

\begin{quote}
\emph{mhaĩ mhārī beṭī parṇāī.}
\end{quote}

\begin{quote}
I gave my daughter in marriage.
\end{quote}

Past (I …): \mw{म्हैं} \mw{बेटो} \mw{परणायो}, \emph{mhaĩ beṭo parṇāyo}.

\begin{cousinweb}
\textbf{\mw{परणावणो}} (\emph{parṇāvṇo}) is inherited from Sanskrit \emph{pariṇāpayati}, \textquotedblleft{}causes a man to marry\textquotedblright{}; CDIAL lists Old Marwari \emph{paraṇāi}, \textquotedblleft{}to give in marriage\textquotedblright{}, under it (CDIAL 7821).
\end{cousinweb}

\begin{grammarlens}[title={What you've built}]
One more word for saying what happened.
\end{grammarlens}

\subsection*{Your turn}
\begin{itemize}
\raggedright
  \item \emph{Say it:} \emph{parṇāvṇo}
  \item \emph{Say it:} \emph{parṇāvṇo}, once more
  \item \emph{Say it:} \emph{mhaĩ mhārī beṭī parṇāī.} — I gave my daughter in marriage.
  \item \emph{Recall:} say the Marwadi for a wedding pavilion, then the Marwadi for henna, then say \emph{parṇāvṇo} again
\end{itemize}

\begin{tcolorbox}[breakable,colback=teal!4,colframe=teal!35!black,title={Before you move on}]
What does \mw{परणावणो} mean? (\textquotedblleft{}To give in marriage; to marry off\textquotedblright{}.) Say it once more.
\end{tcolorbox}
\section[\texorpdfstring{\mw{दायजो} (dāyjo)}{दायजो (dāyjo)}]{\mw{दायजो} (dāyjo) — a dowry}

\label{lesson:MW-C248-dayjo}



\begin{quote}
Before the new one: say the Marwadi for henna, then the Marwadi for to marry off.
\end{quote}

\subsection*{You'll want to know: \mw{दायजो}}
\textbf{\mw{दायजो}} — \emph{dāyjo} — \textquotedblleft{}a dowry\textquotedblright{}.

\begin{quote}
\mw{म्हैं} \mw{दायजो} \mw{कोनी} \mw{लियो।}
\end{quote}

\begin{quote}
\emph{mhaĩ dāyjo konī liyo.}
\end{quote}

\begin{quote}
I took no dowry.
\end{quote}

A masculine noun.

\begin{cousinweb}
\textbf{\mw{दायजो}} (\emph{dāyjo}) is inherited from Sanskrit \emph{dāyādya}, \textquotedblleft{}an inheritance\textquotedblright{}; CDIAL lists Marwari \emph{dāyajo}, \textquotedblleft{}dowry\textquotedblright{}, under it (CDIAL 6290). A dowry is probably named as the bride's share of the family's goods.
\end{cousinweb}

\begin{grammarlens}[title={What you've built}]
One more word for saying what happened.
\end{grammarlens}

\subsection*{Your turn}
\begin{itemize}
\raggedright
  \item \emph{Say it:} \emph{dāyjo}
  \item \emph{Say it:} \emph{dāyjo}, once more
  \item \emph{Say it:} \emph{mhaĩ dāyjo konī liyo.} — I took no dowry.
  \item \emph{Recall:} say the Marwadi for henna, then the Marwadi for to marry off, then say \emph{dāyjo} again
\end{itemize}

\begin{tcolorbox}[breakable,colback=teal!4,colframe=teal!35!black,title={Before you move on}]
What does \mw{दायजो} mean? (\textquotedblleft{}A dowry\textquotedblright{}.) Say it once more.
\end{tcolorbox}
\section[\texorpdfstring{\mw{सजणो} (sajṇo)}{सजणो (sajṇo)}]{\mw{सजणो} (sajṇo) — to dress up; to be adorned}

\label{lesson:MW-C248-sajno}



\begin{quote}
Before the new one: say the Marwadi for to marry off, then the Marwadi for a dowry.
\end{quote}

\subsection*{You'll want to know: \mw{सजणो}}
\textbf{\mw{सजणो}} — \emph{sajṇo} — \textquotedblleft{}to dress up; to be adorned\textquotedblright{}.

\begin{quote}
\mw{ब्याव} \mw{मांय} \mw{सगळा} \mw{सजै} \mw{है।}
\end{quote}

\begin{quote}
\emph{byāv mā̃y sagḷā sajai hai.}
\end{quote}

\begin{quote}
At a wedding everyone dresses up.
\end{quote}

Past (he …): \mw{वो} \mw{छोरो} \mw{सज्यो}, \emph{vo chhoro sajyo}.

\begin{cousinweb}
\textbf{\mw{सजणो}} (\emph{sajṇo}) is probably inherited from Sanskrit \emph{sajjyate}, \textquotedblleft{}is made ready\textquotedblright{}; for a bride or a house, being made ready is being adorned. CDIAL gives no Marwari form in this entry; it lists Hindi \emph{sājnā}, \textquotedblleft{}to be prepared\textquotedblright{} (CDIAL 13093).
\end{cousinweb}

\begin{grammarlens}[title={What you've built}]
One more word for saying what happened.
\end{grammarlens}

\subsection*{Your turn}
\begin{itemize}
\raggedright
  \item \emph{Say it:} \emph{sajṇo}
  \item \emph{Say it:} \emph{sajṇo}, once more
  \item \emph{Say it:} \emph{byāv mā̃y sagḷā sajai hai.} — At a wedding everyone dresses up.
  \item \emph{Recall:} say the Marwadi for to marry off, then the Marwadi for a dowry, then say \emph{sajṇo} again
\end{itemize}

\begin{tcolorbox}[breakable,colback=teal!4,colframe=teal!35!black,title={Before you move on}]
What does \mw{सजणो} mean? (\textquotedblleft{}To dress up; to be adorned\textquotedblright{}.) Say it once more.
\end{tcolorbox}
\section[\texorpdfstring{\mw{सेहरो} (sehro)}{सेहरो (sehro)}]{\mw{सेहरो} (sehro) — a bridegroom's garland}

\label{lesson:MW-C248-sehro}



\begin{quote}
Before the new one: say the Marwadi for a dowry, then the Marwadi for to dress up.
\end{quote}

\subsection*{You'll want to know: \mw{सेहरो}}
\textbf{\mw{सेहरो}} — \emph{sehro} — \textquotedblleft{}a bridegroom's garland\textquotedblright{}.

\begin{quote}
\mw{सेहरो} \mw{पीळो} \mw{अर} \mw{लाल} \mw{है।}
\end{quote}

\begin{quote}
\emph{sehro pīḷo ar lāl hai.}
\end{quote}

\begin{quote}
The bridegroom's garland is yellow and red.
\end{quote}

A masculine noun.

\begin{cousinweb}
\textbf{\mw{सेहरो}} (\emph{sehro}) is inherited from Sanskrit \emph{śekhara}, \textquotedblleft{}a chaplet\textquotedblright{}. CDIAL gives no Marwari form in this entry; it lists Hindi \emph{sehrā}, \textquotedblleft{}a garland worn at a wedding\textquotedblright{} (CDIAL 12604).
\end{cousinweb}

\begin{grammarlens}[title={What you've built}]
One more word for saying what happened.
\end{grammarlens}

\subsection*{Your turn}
\begin{itemize}
\raggedright
  \item \emph{Say it:} \emph{sehro}
  \item \emph{Say it:} \emph{sehro}, once more
  \item \emph{Say it:} \emph{sehro pīḷo ar lāl hai.} — The bridegroom's garland is yellow and red.
  \item \emph{Recall:} say the Marwadi for a dowry, then the Marwadi for to dress up, then say \emph{sehro} again
\end{itemize}

\begin{tcolorbox}[breakable,colback=teal!4,colframe=teal!35!black,title={Before you move on}]
What does \mw{सेहरो} mean? (\textquotedblleft{}A bridegroom's garland\textquotedblright{}.) Say it once more.
\end{tcolorbox}
